% your LaTeX will extract the file if required

\RequirePackage{fix-cm}

\documentclass[twocolumn]{svClass}          % twocolumn
%
\smartqed  % flush right qed marks, e.g. at end of proof
%
\usepackage{graphicx}

% Insert the name of "your journal" with
 \journalname{International Journal of Engineering Insights: }
%
\begin{document}

\title{Title: Preparation of articles for the INJEI journal}
\subtitle{Template}

\author{First Author name \and  Second Author name \and  Third Author name}

\institute{First Author name [ORCID] \at
			  First Institute \\
              San Juan, Argentina \\
              Tel.: +123-45-678910\\ 
              \email{fauthor@example.com}           %  \\
%             \emph{Present address:} of F. Author  %  if needed
           \and
            First Author name \and Second Author name [ORCID] \at
            Second Institute \\
            Hanover, Germany
            \and
            Third Author name [ORCID] \at
            Third Institute \\
            Brisbane, Australia            
}

\date{Received: date / Accepted: date}
% The correct dates will be entered by the editor


\maketitle

\begin{abstract}
These instructions constitute a guide for the preparation of articles for the IJEI Journal. Use this document as a set of instructions. Please use this document as a "template" to prepare your manuscript. For shipping guidelines, follow the instructions of the article submission system on the journal website. The title should not have more than 15 words. The number of authors should not be greater than six. The abstract should not occupy more than 200 words (approximately 10 to 12 lines). Replace Surname1, Initial1. by the last name of the first author and the initial of his name, Surname2, Initial2 by that of the second, etc. Remember that the summary cannot exceed the first page. In the header of page 2 and following should appear the last name of the first author and his initial followed by "et al.".
\keywords{First keyword \and Second keyword \and (5-10 keywords in English )}

\end{abstract}

\section{Introduction}
\label{intro}

This document is a template for Microsoft Word and Latex versions. The language handled by the journal is entirely in English (preferably American).
The first paragraph is not indented; the rest of the paragraphs are indented. 


\section{Figures, Equations, and so on}
\label{sec:1}

\subsection{Figures}
Figures can use all the space (Fig. \ref{fig1}) or the required space (Fig. \ref{fig2}).

\subsection{Equations}
\label{sec:2}
When required. Don't forget to give each section and subsection a unique label (see Sect.~\ref{sec:1}).
\paragraph{Proof:} Use paragraph headings as needed. All numbered equations must be invoked (\ref{eqIKsol}).

\begin{equation} 
	\eta =  \mathbf{J^\dagger}(\mathbf{q})\mathbf{\dot x}, \\
	\label{eqIKsol}
\end{equation}

% For one-column wide figures use
\begin{figure}
  \includegraphics[width=0.48\textwidth]{Figures/fig1}
\caption{Please write your figure caption here}
\label{fig1}       % Give a unique label
\end{figure}

% For two-column wide figures use
\begin{figure*}
  \includegraphics[width=\textwidth]{Figures/fig2}
\caption{Please write your figure caption here}
\label{fig2}       % Give a unique label
\end{figure*}

Do not forget to invoke the Tables, for example: Table \ref{tab1}.
% For tables use
\begin{table}
% table caption is above the table
\caption{Please write your table caption here}
\label{tab1}       % Give a unique label
% For LaTeX tables use
\begin{tabular}{lll}
\hline\noalign{\smallskip}
first & second & third  \\
\noalign{\smallskip}\hline\noalign{\smallskip}
number11 & number12 & number13 \\
number21 & number22 & number23 \\
\noalign{\smallskip}\hline
\end{tabular}
\end{table}


\section{How to invoke References}
All References need be cited in text. When they are, they appear on the line, in square brackets, inside the punctuation, example \cite{palomaki2017wind}. Multiple references can be cited in this way \cite{richardson2019spray,zhang2019application,khurshid2016effects}. 
When citing a section in a book, please give the relevant page numbers. In text, refer simply to the reference number. 

\section{Conclusions}
Conclusions may review the most important points of a scientific article, but should not replicate the abstract in the conclusions. Conclusions may discuss the importance of the work done or suggest applications or future work.


\begin{acknowledgements}
If you'd like to thank anyone, place your comments here
and remove the percent signs.
\end{acknowledgements}


 \section*{Conflict of interest}
The authors declare that they have no conflict of interest.


% BibTeX users please use one of
%\bibliographystyle{spbasic}      % basic style, author-year citations
%\bibliographystyle{spmpsci}      % mathematics and physical sciences
%\bibliographystyle{spphys}       % APS-like style for physics
%\bibliography{}   % name your BibTeX data base

\label{refer}
\bibliographystyle{IEEEtran}
\bibliography{./references}

\end{document}


